\documentclass[12pt,a4paper,oneside,dvipsnames,table,svgnames,skins,theorems]{report}

\usepackage{ProfCollege}
%\usepackage[nominted]{ProfLycee}
\usepackage[tp]{tmbase} %pour le moment placé dans la racine du .tex ... !




\ordi{fixe} %pour fichier image
\matierechap{m} %m pour maths, pc pour physique chimie, rien pour vierge
\version{1} %1 : PROF   2: ELEVE sans corrigé  3: corrigé des exos sans le cours
\setcounter{chapter}{0} %compteur de chapitre

\ifportable

\graphicspath{{C:/Users/Thomas/latex nas/images/}} %chemin du fichier image
\fi
\iffixe
\graphicspath{{C:/Users/Thom/Documents/latex nas/images/}} %chemin du fichier image
\fi

%\author{TP}
\author{Devoir type CCF} %mettre ici non pas l'auteur mais la partie en haut à gauche de l'encadré


%\newcolumntype{C}{>{\centering\arraybackslash}X}
%\renewcommand{\tabularxcolumn}[1]{>{\arraybackslash}m{#1}}
\begin{document}

\setcounter{chapter}{2}

\chapter{ Suites  }





\Opensolutionfile{mycor}[ficcorex]

\vspace{-1cm}
\section*{Évolution salariale}

Le patron d'une grosse entreprise cherche une option pour réduire son salaire pour montrer à ses employés qu'il se soucie de la santé financière et de leur avis. Il leur propose deux options :

\begin{itemize}
\item A : Se verser un salaire symbolique de 1\euro{} qui doublera chaque mois.
\item B : Se verser le même salaire que le plus bas salaire et s'augmenter de 50\euro{} par mois.
\end{itemize}





\problematique[]{Comment déterminer la meilleur proposition pour chaque parti avec une durée fixée à 15 mois ?}

%\newpage

\begin{enumerate}
\item \textbf{Indiquer} le salaire de l'option A au bout de 2 mois.

\Lignespointilles{1}

\item \textbf{Indiquer} le salaire de l'option B au bout de 2 mois.

\Lignespointilles{1}

\item \textbf{Proposer} une méthode pour répondre à la problématique. 

\Lignespointilles{4}

\vspace{0.5cm}

\item \appel[explication de votre méthode.] 
\vspace{0.5cm}

\item \textbf{Mettre en oeuvre} votre méthode.
\Lignespointilles{8}




\vspace{0.5cm}

\item \textbf{Répondre} à la problématique  en ayant un regard critique et en adoptant les différents points de vue. 

\Lignespointilles{2}

\item \textbf{Déterminer} pour chaque cas le nombre de mois nécessaire pour obtenir un salaire de 100000\euro{}
\Lignespointilles{2}
\end{enumerate}





\end{document}